\documentclass{article}
\usepackage{graphicx} % Required for inserting images
\usepackage{amsmath}
\title{Abstract Algebra}
\author{Maxwel Kim}
\date{June 2026}

\begin{document}

\maketitle

\section*{Introduction}
\text{I've understood abstract algebra as 3 branches: \textbf{group},\textbf{ring}, and \textbf{field}. And therefore, I will organize topics into these three branches.}



\section*{Groups}


\textbf{Group:} any set that satisfies \\

\begin{enumerate}
    \item associative property

    \[
    (ab)c=a(bc)
    \]

    \item has identity element

    \[
    ae=ea=a
    \]

    \item inverse element

    \[
    ab=ba=e
    \]
\end{enumerate}

\textbf{Example:} set 

\[
\{1,-1,i,-i\} 
\]

is a group because $-1$ is inverse element of itself, while $i,-i$ are inverse of $-i,i$

\parskip=24

\textbf{theorem:}
group G's identity is unique

proof:

\[
1. \forall a \in G, ae=a
\]

\[
2.\forall a\in G, e'a=a
\]

then 

if $a=e'$, $e'e=e'$, and if $a=e$, $e'e=e$, therefore $e=e'$

\textbf{theorem:} since in groups, there exists an inverse element, cancellation is valid 

\[
ba=ca \rightarrow b=c
\]

leads to \textbf{inverse element being unique} because of cancellation

given

\[
ab=ba=e
\]

if $b$ and $c$ are inverse of $a$, $ab=e$, and $ac=e$. therefore $b=c$


\textbf{Definition:} Order of a group is number of element of $G$ expressed as $|G|$

\textbf{Definition:} Order of an Element is order of $g$, expressed as $|g|$ where $g$ is element of $G$

\textbf{example:} group $U(15)$ where it's in $\mod 15$,

\[
U(15)=\{1,2,3,4,7,8,11,13,14\}
\]

its order of group is $8$ because $$\{3\}$$ is group of units modulo $15$

\textbf{Definition:} Subgroup is when $\forall G, \exists H $ s.t. $H\subset G$ that satisfies to be another group.

\textbf{test for subgroup:} 

\[
\forall H\leq G, \exists \,a,b \text{ such that} \; ab^{-1}\in H \text{ then} \; \textbf{H is subgroup of G}
\]

\textbf{proof:}

\[
\text{since H's binary operation is same as G's,it works for associative property,} \; \text{ for identity,} \; \exists x\in H \text{ such that } e=x*x^{-1} \text{ replace x with a, and b that exists in H, } a=e, b=h^{-1}, h^{-1}=eh^{-1}=ab^{-1}
\]

\textbf{example:} 



\[
H=\{x\in G\, | \, x^2=e\}
\]

let's say $a,b$ are in $H$, then 

\[
a^2=e, b^2=e
\]

then 

\[
\left(ab^{-1}\right)^2=ab^{-1}ab^{-1}=a^2\left(b^{-1}\right)^2=a^2\left(b^2\right)^{-1}=ee^{-1}=e
\]

therefore by \textbf{one-step subgroup test, H is G's Subgroup}




\textbf{Claim: } \langle a \rangle \text{ is a subgroup.}

\textbf{Proof: }
\text{Since } a \in \langle a \rangle,\; \langle a \rangle \neq \varnothing.
\text{ Let } a^m, a^n \in \langle a \rangle.
\text{ Then}
\[
a^m(a^n)^{-1}
= a^m a^{-n}
= a^{m-n}
\in \langle a \rangle.
\]
\text{Therefore, by the subgroup test, } \langle a \rangle \le G.
\qquad \blacksquare

\textbf{Cyclic Groups:}

\[

\text{If group G can be expressed into } G=\{a^n \, | \, n\in \mathbb{Z}\}, \text{ and } a \text{ exists,} \, a \text{ is generator of } \, G, \text{as well as cyclic group generated by it is called } \, G=<a> 
\]

\textbf{Test for $a^i=a^j$}

\textbf{Let } G \text{ be a group and } a \in G.

\textbf{If } a \text{ has infinite order, then}
\[
i = j \iff a^i = a^j.
\]

\textbf{If } a \text{ has finite order } n, then
\[
\langle a \rangle = \{e, a, a^2, \dots, a^{n-1}\},
\]
and
\[
n \mid (i - j) \iff a^i = a^j.
\]

\documentclass{article}
\usepackage{amsmath, amsthm, amssymb}

\newtheorem{theorem}{Theorem}

\begin{document}

\begin{theorem}[Fundamental Theorem of Cyclic Groups]
Let $G = \langle a \rangle$ be a cyclic group with $|G| = n$. Then:
\begin{enumerate}
    \item Every subgroup of $G$ is cyclic.
    \item The order of any subgroup of $G$ divides $n$.
    \item For each positive divisor $k$ of $n$, there exists exactly one subgroup of order $k$, namely $\langle a^{n/k} \rangle$.
\end{enumerate}
\end{theorem}


\begin{proof}
\textbf{Every subgroup is cyclic:} Let $H \leq G$. If $H = \{e\}$, it is cyclic. Otherwise, 
let $m$ be the smallest positive integer such that $a^m \in H$. We claim $H = \langle a^m \rangle$.

For any $b \in H$, write $b = a^k$. By the division algorithm, $k = qm + r$ with $0 \leq r < m$. Then:
\[
    a^r = a^{k-qm} = a^k \cdot (a^m)^{-q} \in H
\]
Since $m$ is minimal and $0 \leq r < m$, we must have $r = 0$, so $b = (a^m)^q \in \langle a^m \rangle$. 
Hence $H = \langle a^m \rangle$.

\textbf{Uniqueness of subgroup of order $k$:} Let $k \mid n$. Then:
\[
    |\langle a^{n/k} \rangle| = \frac{n}{\gcd(n,\, n/k)} = k
\]
If $H = \langle a^m \rangle$ has order $k$, then:
\[
    |H| = \frac{n}{\gcd(n,m)} = \frac{n}{m} = k \implies m = \frac{n}{k}
\]
Thus $H = \langle a^{n/k} \rangle$ is the unique subgroup of order $k$.
\end{proof}

\documentclass{article}
\usepackage{amsmath, amsthm, amssymb}

\newtheorem{theorem}{Theorem}

\begin{document}

\begin{theorem}[Number of Elements of Each Order in a Cyclic Group]
For each positive divisor $d$ of $n$, the number of elements of order $d$ 
in a cyclic group of order $n$ is $\phi(d)$.
\end{theorem}

\begin{proof}
By Theorem 3.3, the group has exactly one subgroup of order $d$, 
call it $\langle a \rangle$. Every element of order $d$ must generate 
$\langle a \rangle$, and by Theorem 3.2 (Corollary), $a^k$ generates 
$\langle a \rangle$ if and only if $\gcd(k, d) = 1$. Therefore the 
number of elements of order $d$ is $\phi(d)$.
\end{proof}

\noindent\textbf{Remark.} The number of elements of order $d$ in a cyclic 
group of order $n$ depends only on $d$. Thus each of $\mathbb{Z}_8$, 
$\mathbb{Z}_{640}$, and $\mathbb{Z}_{80000}$ contains exactly 
$\phi(8) = 4$ elements of order $8$.

\bigskip
\noindent\textbf{Table of $\phi(n)$:}

\begin{center}
\begin{tabular}{c|cccccccccccc}
$n$ & 1 & 2 & 3 & 4 & 5 & 6 & 7 & 8 & 9 & 10 & 11 & 12 \\
\hline
$\phi(n)$ & 1 & 1 & 2 & 2 & 4 & 2 & 6 & 4 & 6 & 4 & 10 & 4 \\
\end{tabular}
\end{center}

\documentclass{article}
\usepackage{amsmath, amssymb}

\begin{document}

\textbf{Definition (Group Isomorphism).}
An isomorphism $\phi$ from a group $G$ to a group $\bar{G}$ is a bijective
function from $G$ to $\bar{G}$ that preserves the group operation. That is,
for all elements $a, b \in G$:
\[
    \phi(ab) = \phi(a)\phi(b)
\]
If an isomorphism from $G$ to $\bar{G}$ exists, $G$ and $\bar{G}$ are said
to be isomorphic, denoted $G \approx \bar{G}$.

\medskip
\noindent
That isomorphic groups have the same order follows immediately from the
definition. Also note that in the defining equation, the operation on the
left side is that of $G$, while the operation on the right side is that
of $\bar{G}$.


\textbf{Example}
The map $\phi(x) = x^3$ from $\mathbb{R}$ to $\mathbb{R}$ under multiplication is \textit{not} an isomorphism.
Although $\phi$ is a bijection, it does not preserve the operation, since for all $x, y$,
$(x + y)^3 \neq x^3 + y^3$ in general.

\textbf{Example}
$U(10) \not\cong U(12)$.
The proof is slightly involved.
Suppose $\phi$ is an isomorphism from $U(10)$ to $U(12)$.
Every element $x$ of $U(12)$ satisfies $x^2 = 1$, so

\[
\phi(9) = \phi(3 \cdot 3) = \phi(3)\phi(3) = 1
\]

and

\[
\phi(1) = \phi(1 \cdot 1) = \phi(1)\phi(1) = 1.
\]

Therefore $\phi(9) = \phi(1)$.
However, $9 \neq 1$, which contradicts the assumption that $\phi$ is injective (one-to-one).

\begin{definition}[Coset of $H$ in $G$]
Let $H$ be an arbitrary subset of a group $G$ and let $a \in G$. Define the following sets:
\begin{align*}
aH &= \{ ah \mid h \in H \} \\
Ha &= \{ ha \mid h \in H \} \\
aHa^{-1} &= \{ aha^{-1} \mid h \in H \}
\end{align*}
In particular, if $H$ is a subgroup of $G$, then $aH$ is called the \textbf{left coset} of $H$ in $G$, and $Ha$ is called the \textbf{right coset} of $H$ in $G$. The element $a$ is called the 
\end{definition}

\textbf{example:} let $G=S_3$, $H=\{(1),(13)\}$ then

\begin{align*}
(1)H &= H \\
(12)H &= \{(12),(12)(13)\} = \{(12),(132)\} = (132)H \\
(13)H &= \{(13),(1)\} = H \\
(23)H &= \{(23),(23)(13)\} = \{(23),(123)\} = (123)H
\end{align*}

\textbf{example}: let $\mathbb{Z}_9$'s addition group $H=\{0,3,6\}$, then

\begin{align*}
0 + H &= \{(0,3,6\} = 3 + H =6 + H \\
1 + H &= \{1,4,7\} = 4+H = 7+H\\
2+H &= \{2,5,8\} = 5+H=8+H\\
\end{align*}

\textbf{Lagrange's Theorem}\\

For all finite group $G$ and subgroup $H$, $|H|$ is a divisor of $|G|$. Furthermore, $|G|/|H|$ is the number of cosets.

\textbf{proof} \\
Set $H$'s different left coset as $a_1H,a_2H...a_rH$ then there exists $aH=a_iH$ for any $i$, since $a\in aH$, $a\in a_iH$ therefore $G=a_1H\cup a_2H \cup ...\cup a_rH$ and since they are all different, not common multiple, $|G|=|a_1H|+|a_2H|+...+|a_rH|$ therefore $|G|= r|H|$

\begin{theorem}
The order of an element in a finite direct product of groups is the least common multiple of the orders of its components. That is, for $(g_1, g_2, \ldots, g_n) \in G_1 \oplus G_2 \oplus \cdots \oplus G_n$,
\[
|(g_1, g_2, \ldots, g_n)| = \text{lcm}(|g_1|, |g_2|, \ldots, |g_n|)
\]
\end{theorem}

\begin{proof}
Let $s = |(g_1, g_2, \ldots, g_n)|$ and $t = \text{lcm}(|g_1|, |g_2|, \ldots, |g_n|)$.

Since $(g_1, g_2, \ldots, g_n)^t = (g_1^t, g_2^t, \ldots, g_n^t) = (e_1, e_2, \ldots, e_n)$,
we have $s \leq t$.

Conversely, $(g_1, g_2, \ldots, g_n)^s = (e_1, e_2, \ldots, e_n)$ implies $g_i^s = e_i$ for all $i$,
so $|g_i|$ divides $s$ for each $i$. Hence $t \mid s$, giving $t \leq s$.

Therefore $s = t$.
\end{proof}


\begin{theorem}
For finite cyclic groups $G$ and $H$, $G \oplus H$ is cyclic if and only if $\gcd(|G|, |H|) = 1$.
\end{theorem}

\begin{proof}
Let $|G| = m$, $|H| = n$.

$(\Rightarrow)$ Suppose $G \oplus H$ is cyclic with generator $(g, h)$ and $\gcd(m,n) = d$. Then:
\[
(g, h)^{mn/d} = \left((g^m)^{n/d},\, (h^n)^{m/d}\right) = (e, e)
\]
so $|(g,h)| \leq mn/d$. But $|(g,h)| = mn$ since $(g,h)$ generates $G \oplus H$, forcing $d = 1$.

$(\Leftarrow)$ Suppose $\gcd(m, n) = 1$. Let $G = \langle g \rangle$, $H = \langle h \rangle$. Then:
\[
|(g, h)| = \text{lcm}(m, n) = mn = |G \oplus H|
\]
so $(g, h)$ generates $G \oplus H$, making it cyclic.
\end{proof}

\textbf{Definition:} Normal Subgroup is when $aH=Ha$ applies in $H\in G$

\textbf{test for Normal Subgroup:} 

\[
\forall x\in G, xHx^{-1}\subseteq H
\]

\textbf{proof:}\\
this one's easy to prove, if $H$ is normal subgroup of $G$, then for $x\in G$ and $h\in H$, $h'\in H$ exists such that $xh=h'x$ applies. therefore $xhx^{-1}=h'$, $xHx^{-1}\subseteq H$. 

\textbf{example:} every subgroup of abellian group is normal subgroup($ah=ha$ applies as definition)

\textbf{Definition:} Group Homomorphism 

\[
\forall a,b\in G, \text{function } \phi \text{ where it goes } G\rightarrow\overline{G}, \;\phi(ab)=\phi(a)*\phi(b)
\]

And Kernel

\[
\{x\in G|\phi(x)=e\}
\]

\begin{theorem}
Let $\phi$ be a homomorphism from group $G$ to group $\bar{G}$, and let $H$ be a subgroup of $G$. Then the following hold:
\begin{enumerate}
    \item $\phi(H) = \{\phi(h) \mid h \in H\}$ is a subgroup of $\bar{G}$.
    \item If $H$ is cyclic, then $\phi(H)$ is also cyclic.
    \item If $H$ is abelian, then $\phi(H)$ is also abelian.
    \item If $H$ is a normal subgroup of $G$, then $\phi(H)$ is normal in $\phi(G)$.
    \item If $|\mathrm{Ker}\,\phi| = n$, then $\phi$ is an $n$-to-$1$ map from $G$ onto $\phi(G)$.
    \item If $|H| = n$, then $|\phi(H)|$ divides $n$.
    \item If $\bar{K}$ is a subgroup of $\bar{G}$, then $\phi^{-1}(\bar{K}) = \{k \in G \mid \phi(k) \in \bar{K}\}$ is a subgroup of $G$.
    \item If $\bar{K}$ is a normal subgroup of $\bar{G}$, then $\phi^{-1}(\bar{K}) = \{k \in G \mid \phi(k) \in \bar{K}\}$ is a normal subgroup of $G$.
    \item If $\phi$ is surjective and $\mathrm{Ker}\,\phi = \{e\}$, then $\phi$ is an isomorphism from $G$ to $\bar{G}$.
\end{enumerate}
\end{theorem}

\begin{proof}
Since a homomorphism preserves the group operation, properties 4, 3, 2 follow from Theorem 5.3, and properties 1, 2, 3 above correspond to properties 3, 2 of Theorem 5.3, respectively.

To prove property 4: for $\phi(h) \in \phi(H)$ and $\phi(g) \in \phi(G)$, since $H$ is normal in $G$,
\[
\phi(g)\phi(h)\phi(g)^{-1} = \phi(ghg^{-1}) \in \phi(H).
\]

Property 5 follows from property 6 of Theorem 9.1 and the fact that each coset of $\mathrm{Ker}\,\phi = \phi^{-1}(e)$ has the same number of elements. To prove property 6: let $\phi_H$ denote the restriction of $\phi$ to $H$ (i.e., $\phi$ with domain restricted to $H$ and codomain $\phi(H)$). By property 5, $\phi_H$ is a $t$-to-$1$ map, so $|\phi(H)| \cdot t = n$, which proves the property.

Since $e \in \phi^{-1}(\bar{K})$, $\phi^{-1}(\bar{K})$ is nonempty. For any $k_1, k_2 \in \phi^{-1}(\bar{K})$, we have $\phi(k_1), \phi(k_2) \in \bar{K}$, so $\phi(k_2)^{-1} \in \bar{K}$ and $\phi(k_1 k_2^{-1}) = \phi(k_1)\phi(k_2)^{-1} \in \bar{K}$. Therefore $k_1 k_2^{-1} \in \phi^{-1}(\bar{K})$, which by the one-step subgroup test proves property 7.

Property 8 is proved using the normality criterion from Theorem 9.1. Since $\bar{K}$ is normal in $\bar{G}$, for any $x \in G$ and $k \in \phi^{-1}(\bar{K})$, we have $\phi(k) \in \bar{K}$, so
\[
\phi(xkx^{-1}) = \phi(x)\phi(k)(\phi(x))^{-1} \in \bar{K}.
\]
Therefore $xkx^{-1} \in \phi^{-1}(\bar{K})$, proving property 8.

Finally, property 9 follows from property 5.
\end{proof}

\textbf{example:} for $\phi(x)=x^4$, since $(xy)^4=x^4y^4$, $\phi$ is homomorphism, and Ker$\phi=\{x|x^4=1\}=\{1,-1,i,-2\}$

\begin{theorem}
Let $\phi : G \to \bar{G}$ be a homomorphism from group $G$ to group $\bar{G}$.
Then the map
\[
  \Psi : G/\mathrm{Ker}\,\phi \;\longrightarrow\; \phi(G),
  \qquad
  g\,\mathrm{Ker}\,\phi \;\mapsto\; \phi(g)
\]
is an isomorphism from $G/\mathrm{Ker}\,\phi$ to $\phi(G)$. That is,
\[
  G/\mathrm{Ker}\,\phi \;\cong\; \phi(G).
\]
\end{theorem}

\begin{proof}
Call the map $g\,\mathrm{Ker}\,\phi \mapsto \phi(g)$ by the name $\Psi$.
By property~5 of Theorem~10.1, $\Psi$ is well-defined
(i.e.\ its value is independent of the choice of coset representative)
and is injective. Moreover,
\begin{align*}
  \Psi(x\,\mathrm{Ker}\,\phi \cdot y\,\mathrm{Ker}\,\phi)
  &= \Psi(xy\,\mathrm{Ker}\,\phi) \\
  &= \phi(xy) \\
  &= \phi(x)\,\phi(y) \\
  &= \Psi(x\,\mathrm{Ker}\,\phi)\,\Psi(y\,\mathrm{Ker}\,\phi),
\end{align*}
so $\Psi$ preserves the group operation and is therefore a homomorphism.
Since $\Psi$ is both injective and surjective onto $\phi(G)$, it is an isomorphism.
\end{proof}

\documentclass{article}
\usepackage{amsmath, amsthm, amssymb}

\newtheorem{theorem}{Theorem}[section]

\begin{document}

\setcounter{section}{11}

%% ============================================================
%% Theorem 11.3 -- Sylow's First Theorem
%% ============================================================
\begin{theorem}[Existence of Subgroups of Prime-Power Order (Sylow's First Theorem, 1872)]
Let $G$ be a finite group and $p$ a prime. If $p^k$ divides $|G|$, then $G$ has at least one subgroup of order $p^k$.
\end{theorem}

\begin{proof}
We apply mathematical induction on $|G|$. If $|G|=1$, Theorem 11.3 is trivially true. Now assume the theorem holds for all groups of order less than $|G|$.

If $G$ has a proper subgroup $H$ such that $p^k$ divides $|H|$, then by the induction hypothesis $H$ contains a subgroup of order $p^k$, and the theorem is proved.

Therefore, assume that no proper subgroup of $G$ has order divisible by $p^k$.

Consider the class equation of $G$:
\[
|G| = |Z(G)| + \sum |G : C(a)|,
\]
where the sum is taken over one representative $a$ from each conjugacy class $\mathrm{cl}(a)$ with $a \notin Z(G)$.

Since $p^k$ divides $|G| = |G:C(a)|\cdot|C(a)|$ and $p^k$ does not divide $|C(a)|$, we have that $p$ divides $|G:C(a)|$ for every $a \notin Z(G)$. From the class equation it follows that $p$ divides $|Z(G)|$.

By the Fundamental Theorem of Finite Abelian Groups (Theorem 10.1), or by Theorem 9.5, $Z(G)$ has an element $x$ of order $p$. Since $x$ is a central element of $G$, $\langle x \rangle$ is a normal subgroup of $G$ and we can form the quotient group $G/\langle x \rangle$.

Now $p^{k-1}$ divides $|G/\langle x \rangle|$, so by the induction hypothesis $G/\langle x \rangle$ has a subgroup of order $p^{k-1}$. By Exercise 21 of Chapter 9, this subgroup has the form $H/\langle x \rangle$ for some subgroup $H$ of $G$.

Finally, since $|H/\langle x \rangle| = p^{k-1}$ and $|\langle x \rangle| = p$, we get $|H| = p^k$. This completes the proof.
\end{proof}

%% ============================================================
%% Theorem 11.4 -- Sylow's Second Theorem
%% ============================================================
\begin{theorem}[Sylow's Second Theorem]
If $H$ is a subgroup of a finite group $G$ and $|H|$ is a power of a prime $p$, then $H$ is contained in some Sylow $p$-subgroup of $G$.
\end{theorem}

\begin{proof}
Let $K$ be a Sylow $p$-subgroup of $G$, and let $\mathcal{C} = \{K_1, K_2, \ldots, K_n\}$ be the set of all conjugates of $K$ in $G$, where $K = K_1$.

Let $S_{\mathcal{C}}$ be the set of all permutations of $\mathcal{C}$. For each $g \in G$, define $\phi_g : \mathcal{C} \to \mathcal{C}$ by
\[
\phi_g(K_i) = g K_i g^{-1}.
\]
One easily verifies that $\phi_g \in S_{\mathcal{C}}$.

Define a map $T : G \to S_{\mathcal{C}}$ by $T(g) = \phi_g$. Since
\[
\phi_{gh}(K_i) = (gh)K_i(gh)^{-1} = g(hK_i h^{-1})g^{-1} = \phi_g(\phi_h(K_i)) = (\phi_g \circ \phi_h)(K_i),
\]
we have $\phi_{gh} = \phi_g \circ \phi_h$, so $T$ is a homomorphism from $G$ to $S_{\mathcal{C}}$.

Consider the image $T(H)$ of $H$ under $T$. Since $|H|$ is a power of $p$, so is $|T(H)|$ (see property 6 of Theorem 9.2). By the Orbit-Stabilizer Theorem (Theorem 6.3), $|\mathrm{orb}_{T(H)}(K_i)|$ divides $|T(H)|$ for each $i$, so $|\mathrm{orb}_{T(H)}(K_i)|$ is a power of $p$.

Now consider when $|\mathrm{orb}_{T(H)}(K_i)| = 1$. This means $\phi_g(K_i) = gK_ig^{-1} = K_i$ for all $g \in H$, i.e., $H \leq N(K_i)$.

The elements of $N(K_i)$ whose order is a power of $p$ are exactly the elements of $K_i$ (see Exercise 5). Therefore, the necessary and sufficient condition for $|\mathrm{orb}_{T(H)}(K_i)| = 1$ is $H \leq K_i$.

To complete the proof, we need to show that $|\mathrm{orb}_{T(H)}(K_i)| = 1$ for some $i$. By Theorem 11.1, $|\mathcal{C}| = |G : N(K)|$ (see Exercise 11). Since $|G:K| = |G:N(K)|\cdot|N(K):K|$ is not divisible by $p$, $|\mathcal{C}|$ is also not divisible by $p$.

Since the orbits partition $\mathcal{C}$, $|\mathcal{C}|$ is the sum of orbit sizes, all of which are powers of $p$. If no orbit had size 1, then $p$ would divide $|\mathcal{C}|$, a contradiction. Therefore an orbit of size 1 exists, and the proof is complete.
\end{proof}

%% ============================================================
%% Theorem 11.5 -- Sylow's Third Theorem
%% ============================================================
\begin{theorem}[Sylow's Third Theorem]
Let $p$ be a prime and $G$ a group of order $p^n m$, where $p$ does not divide $m$. Then the number $n_p$ of Sylow $p$-subgroups of $G$ satisfies:
\begin{enumerate}
    \item $n_p \equiv 1 \pmod{p}$,
    \item $n_p$ divides $m$,
    \item any two Sylow $p$-subgroups of $G$ are conjugate.
\end{enumerate}
\end{theorem}

\begin{proof}
Let $K$ be a Sylow $p$-subgroup of $G$ and let $\mathcal{C} = \{K_1, K_2, \ldots, K_n\}$ be the set of all conjugates of $K$ in $G$, where $K = K_1$.

\textbf{Step 1: $n \equiv 1 \pmod{p}$.}
Define $S_{\mathcal{C}}$, $T$, and the action of $K$ on $\mathcal{C}$ via $T$ exactly as in the proof of Theorem 11.4, but now consider the image $T(K)$ instead of $T(H)$.

As before, $|\mathrm{orb}_{T(K)}(K_i)|$ is a power of $p$, and $|\mathrm{orb}_{T(K)}(K_i)| = 1$ if and only if $K \leq K_i$. Since $K$ and $K_i$ are both Sylow $p$-subgroups (hence of equal order), $K \leq K_i$ implies $K = K_i$. Therefore the unique orbit of size 1 is $\{K_1\} = \{K\}$, and every other orbit has size a multiple of $p$ greater than 1.

Since the orbits partition $\mathcal{C}$, we get $n = |\mathcal{C}| \equiv 1 \pmod{p}$.

\textbf{Step 2: Every Sylow $p$-subgroup belongs to $\mathcal{C}$.}
Suppose $H$ is a Sylow $p$-subgroup not in $\mathcal{C}$. Using $T(H)$ as before, $|\mathcal{C}|$ is the sum of orbit sizes under $T(H)$, each a power of $p$. Since $H \notin \mathcal{C}$, no orbit has size 1, so $p$ divides $|\mathcal{C}|$. But from Step 1, $|\mathcal{C}| \equiv 1 \pmod{p}$, a contradiction. Hence $H \in \mathcal{C}$, and $\mathcal{C}$ is exactly the set of all Sylow $p$-subgroups, so $n_p = |\mathcal{C}|$.

\textbf{Step 3: $n_p$ divides $m$.}
This follows directly from $n_p = |\mathcal{C}| = |G : N(K)|$ (see Exercise 11).Í
\end{proof}



\section*{Ring}

\textbf{Ring} is when this applies:

\[
\begin{aligned}
&1.\ a+b=b+a \\
&2.\ (a+b)+c=a+(b+c) \\
&3.\ \text{identity in respect to addition} \\
&4.\ \forall\, a \in R,\ \exists\, {-a} \in R,\ a+(-a)=0 \\
&5.\ a(bc)=(ab)c \\
&6.\ a(b+c)=ab+ac,\quad (b+c)a=ba+ca
\end{aligned}
\]

\textbf{Example:} 

\[
\mathbb{Z} \text{ is both addition and multiplication wise identity existing, and } \{1,-1\}
\]


\textbf{Test for Subring} if 

\[
a-b, ab\in S \forall a,b
\]

they are subring.

\textbf{Ideal and factor rings:}

it it called ideal when  

\[
\forall r\in  R, a\in A \exists ra\in A , ar\in A
\]

\textbf{Test for Ideal}
\begin{align*}
&1.\ \forall\, a,b \in A,\ a-b \in A \\
&2.\ \forall\, a \in A\ \text{and}\ r \in R,\ ra \in A\ \text{and}\ ar \in A
\end{align*}

\textbf{Example:} for 

\[
n\mathbb{Z}=\{0,\pm n,\pm2n,...\} 
\]

is $\mathbb{Z}$'s Ideal

\textbf{Factor Ring's existence}

if R is ring and A is R's subring, then $\{r+a|r\in R\}$ being $(s+A)+(t+A)=(s+t)+A,\, (s+A)(t+A)=st+A$ is equivalent for $A$ being $R$'s Ideal

\textbf{example:} 

\[
Z/4Z=\{0+4Z,1+4Z,2+4Z,3+4Z\}
\]

then for two elemtns, 

\[
(2+4Z)+(3+4Z)=5+4Z=1=4+4Z=1+4Z
\]
\[
(2+4Z)(3+4Z)=6+4Z=2+4+4Z=2+4Z
\]

\textbf{Definition}

it is called prime ideal when 

\[
\forall a,b\in R, \text{ if } ab\in A, \text{ it's } a\in A \text{ or } b\in A
\]

it is called maximal ideal when 

\[
\text{if } A\subseteq B \subseteq R \text{ then } B=A \text{ or } B=R 
\]

\textbf{Definition:} It is called ring homomorphism when 

\[
\phi(a+b)=\phi(a)+\phi(b), \phi(ab)\phi(a)\phi(b)\]
\]

When it's injective and bijective, it is called ring isomorphism.

\textbf{example:} 

$\phi$ is ring homomorphism on behalf of $\mathbb{Z}_4 \rightarrow \mathbb{Z}_{10}$ on $x\rightarrow 5x$ because 

\[
\phi(x+y)=\phi(x)+\phi(y), 5(x+y)=5x+5y
\]

and 

\[
5xy=5x*5y
\]

on $\mathbb{Z}_{10}$

\begin{theorem}[Properties of Ring Homomorphisms]
Let $\phi$ be a ring homomorphism from a ring $R$ to a ring $S$. Let $A$ be a subring of $R$ and $B$ be an ideal of $S$.

\begin{enumerate}
    \item For any $r \in R$ and any positive integer $n$,
    \[
        \phi(nr) = n\phi(r) \quad \text{and} \quad \phi(r^n) = (\phi(r))^n.
    \]

    \item $\phi(A) = \{\phi(a) \mid a \in A\}$ is a subring of $S$.

    \item If $A$ is an ideal and $\phi$ is onto $S$, then $\phi(A)$ is an ideal.

    \item $\phi^{-1}(B) = \{r \in R \mid \phi(r) \in B\}$ is an ideal of $R$.

    \item If $R$ is commutative, then $\phi(R)$ is commutative.

    \item If $R$ has unity $1$ and $S \neq \{0\}$ and $\phi$ is onto, then $\phi(1)$ is the unity of $S$.

    \item $\phi$ is an isomorphism if and only if $\phi$ is onto and
    \[
        \ker\phi = \{r \in R \mid \phi(r) = 0\} = \{0\}.
    \]

    \item If $\phi$ is an isomorphism from $R$ to $S$, then $\phi^{-1}$ is an isomorphism from $S$ to $R$.
\end{enumerate}
\end{theorem}

\begin{proof}
The proofs follow by the same methods as in Theorems 9.1 and 9.2.
\end{proof}


\textbf{Polynomial Rings}

if $R$ is abelian Ring, then 

\[
R[x] = \{a_nx^n+a_{n-1}x^{n-1}+...+a_0|a_i\in R\}
\]
 
 is  called polynomials in terms of $x$

\begin{theorem}[$D$ an Integral Domain Implies $D[x]$ is an Integral Domain]
If $D$ is an integral domain, then $D[x]$ is an integral domain.
\end{theorem}

\begin{proof}
We already know $D[x]$ is a commutative ring with unity and no zero-divisors. Let
\[
f(x) = a_n x^n + a_{n-1}x^{n-1} + \cdots + a_1 x + a_0, \quad
g(x) = b_m x^m + b_{m-1}x^{m-1} + \cdots + b_1 x + b_0,
\]
with $a_n \neq 0$, $b_m \neq 0$. The leading term of $f(x)g(x)$ is $a_n b_m$, and since $D$ is an integral domain, $a_n b_m \neq 0$.
\end{proof}

\begin{theorem}[Division Algorithm for $F[x]$]
Let $F$ be a field and $f(x), g(x) \in F[x]$ with $g(x) \neq 0$. Then there exist unique
$q(x), r(x) \in F[x]$ such that
\[
f(x) = g(x)q(x) + r(x),
\]
where $r(x) = 0$ or $\deg r(x) < \deg g(x)$.
\end{theorem}

\begin{proof}
By the Well-Ordering Principle applied to $\deg f(x) - \deg g(x)$, we induct on $\deg f(x)$.
If $f(x) = 0$ or $\deg f(x) < \deg g(x)$, take $q(x) = 0$ and $r(x) = f(x)$.

Otherwise, let $\deg f(x) = n$, $\deg g(x) = m$, and write
\[
f(x) = a_n x^n + \cdots + a_0, \quad g(x) = b_m x^m + \cdots + b_0.
\]
Set $f_1(x) = f(x) - a_n b_m^{-1} x^{n-m} g(x)$, so $\deg f_1(x) < \deg f(x)$.
By induction, $f_1(x) = g(x)q_1(x) + r(x)$, giving
\[
f(x) = g(x)\bigl(a_n b_m^{-1} x^{n-m} + q_1(x)\bigr) + r(x).
\]
Uniqueness follows from the fact that if two such representations existed, their difference would yield a contradiction on degrees.
\end{proof}

\begin{example}
Divide $f(x) = 3x^4 + x^3 + 2x^2 + 1$ by $g(x) = x^2 + 4x + 2$ in $\mathbb{Z}_5[x]$:
\[
x^2 + 4x + 2 \;\Big|\; 3x^4 + x^3 + 2x^2 + 1 \;\longrightarrow\; q(x) = 3x^2 + 4x, \quad r(x) = 2x+1.
\]
Thus $3x^4 + x^3 + 2x^2 + 1 = (x^2 + 4x + 2)(3x^2 + 4x) + (2x+1)$.
\end{example}


\textbf{ Definition:} it is called Principal Ideal Domain, PID when the Ideal is in form of 

\[
<a>=\{ra|r\in R\}
\]

\textbf{Theorem} if $F$ is a field, $F[x]$ is a PID \\


\textbf{proof}
as the theorem states, $F[x]$ is Ideal Domain.\\
say $I$ is Ideal of $F[x]$, then $g(x)\in I$, $<g(x)>\subseteq I$, then $\deg r(x)<\deg g(x)$ so $f(x)\in <g(x)>$, $I\subseteq<g(x)>$

\begin{definition}[Irreducible / Reducible Polynomial]
Let $D$ be an integral domain. A non-zero, non-unit $f(x) \in D[x]$ is \textbf{irreducible} over $D$ if whenever $f(x) = g(x)h(x)$ with $g(x), h(x) \in D[x]$, then $g(x)$ or $h(x)$ is a unit in $D[x]$. Otherwise, $f(x)$ is \textbf{reducible}.

In particular, for a field $F$, a non-constant $f(x) \in F[x]$ is irreducible iff $f(x)$ cannot be expressed as a product of two polynomials of lower degree.
\end{definition}

\begin{example} \textbf{example:}
$f(x) = 2x^2 + 4$ is irreducible over $\mathbb{Q}$, but reducible over $\mathbb{Z}$, since
\[
2x^2 + 4 = 2(x^2 + 2),
\]
where neither $2$ nor $x^2 + 2$ is a unit in $\mathbb{Z}[x]$.
\end{example}
 

\textbf{theorem}

Let $F$ be a field and $f(x) \in F[x]$ being $\deg f(x) = 2, 3$ then $f(x)$ being reducible in $F$ is equivalent as $f(x)$ having root in $F$

\textbf{proof:}

if $f(x)=g(x)h(x)$, then $\deg f(x)=\deg g(x) + \deg h(x)$, and since they are 2 or 3, either $g(x)$ or $h(x)$ will be $1$-degree polynomial. This also means if $g(x)=ax+b$, then $g(x)$'s root is $-a^{-1}b$, also being $f(x)$'s root. inversely, if $f(a)=0$ as $a\in F$, as factorization theorem, $x-a$ will be $f(x)$'s factor. therefore $f(x)$ is $F$'s reducible polynomial

\begin{theorem}[Eisenstein's Criterion]
Let $f(x) = a_nx^n + \cdots + a_0 \in \mathbb{Z}[x]$. If a prime $p$ satisfies
\[
p \nmid a_n, \quad p \mid a_{n-1}, \dots, p \mid a_0, \quad p^2 \nmid a_0,
\]
then $f(x)$ is irreducible over $\mathbb{Q}$.
\end{theorem}

\begin{proof}
Suppose $f(x) = g(x)h(x)$ with $g(x) = b_rx^r + \cdots + b_0$ and $h(x) = c_sx^s + \cdots + c_0$.
Since $p \mid a_0 = b_0c_0$ but $p^2 \nmid a_0$, say $p \mid b_0$ and $p \nmid c_0$.
Since $p \nmid a_n = b_rc_s$, we have $p \nmid b_r$, so let $t$ be the smallest index with $p \nmid b_t$.
From $a_t = b_tc_0 + b_{t-1}c_1 + \cdots + b_0c_t$, every term but $b_tc_0$ is divisible by $p$, so $p \nmid a_t$. But $t < n$, contradicting $p \mid a_t$.
\end{proof}

\begin{corollary}[$p$-th Cyclotomic Polynomial]
For any prime $p$,
\[
\Phi_p(x) = x^{p-1} + x^{p-2} + \cdots + 1
\]
is irreducible over $\mathbb{Q}$.
\end{corollary}

\begin{proof}
Set $f(x) = \Phi_p(x+1) = x^{p-1} + \binom{p}{1}x^{p-2} + \cdots + \binom{p}{p-1}$.
All coefficients except the leading one are divisible by $p$, and the constant term $p$ satisfies $p^2 \nmid p$. By Eisenstein, $f(x)$ is irreducible, hence so is $\Phi_p(x)$.
\end{proof}


\begin{theorem}
$\langle p(x) \rangle$ is maximal in $F[x]$ iff $p(x)$ is irreducible over $F$.
\end{theorem}

\begin{proof}
$(\Rightarrow)$ Suppose $\langle p(x)\rangle$ is maximal. If $p(x)=g(x)h(x)$, then
$\langle p(x)\rangle \subseteq \langle g(x)\rangle \subseteq F[x]$, so either $\langle g(x)\rangle = \langle p(x)\rangle$ or $\langle g(x)\rangle = F[x]$, forcing $\deg g(x) = \deg p(x)$ or $\deg g(x)=0$. So $p(x)$ is irreducible.

$(\Leftarrow)$ Suppose $p(x)$ is irreducible and $\langle p(x)\rangle \subseteq I \subseteq F[x]$. Since $F[x]$ is a PID, $I=\langle g(x)\rangle$ and $p(x)=g(x)h(x)$. By irreducibility, $g(x)$ or $h(x)$ is constant, giving $I=F[x]$ or $I=\langle p(x)\rangle$. Thus $\langle p(x)\rangle$ is maximal.
\end{proof}

\begin{corollary}
If $p(x)$ is irreducible over $F$, then $F[x]/\langle p(x)\rangle$ is a field.
\end{corollary}

\textbf{extension field:} when

\[
F\subseteq E 
\]

that $E$'s operation is limited on $F$'s, it's called as \textbf{$E$ is $F$'s Extension Field}

\textbf{example:} given 

\[
f(x)=x^5+2x^2+2x+2\in \mathbb{Z}_3[x]
\]

then factorization is

\[
(x^2+1)(x^3+2x+2)
\]

using this fact, 

\[
E=\mathbb{Z}_3[x]/<x^2+1>
\]

number of element is $9$, or 

\[
E=\mathbb{Z}_3[x]/<x^3+2x+2>
\]

has element of 27

\textbf{Definition: Splitting field} when $f(x)$ is factorized by $E[x]$ as $1$st order, it's called splitting field

\textbf{example:} given

\[
f(x)=x^2+1\in Q[x], 
\]

this factorizes to 

\[
x^2+1=(x+\sqrt{-1})(x-\sqrt{-1})
\]

and therefore this splits in $C$

\documentclass{amsart}
\usepackage{amsmath, amsthm, amssymb}

\newtheorem{theorem}{Theorem}
\newtheorem{remark}{Remark}

\begin{document}

\begin{theorem}[Fundamental Theorem of Galois Theory]
Let $F$ be a field of characteristic $0$ or finite, $E/F$ a splitting field of some $f(x)\in F[x]$,
$G = \mathrm{Gal}(E/F)$. For $K$ with $F\subseteq K\subseteq E$, define $\phi:\mathcal{F}\to\mathcal{G}$
by $\phi(K)=\mathrm{Gal}(E/K)$. Then:
\begin{enumerate}
  \item $[E:K]=|\mathrm{Gal}(E/K)|$ and $[K:F]=|G|/|\mathrm{Gal}(E/K)|$.
  \item $K/F$ Galois $\Rightarrow$ $\mathrm{Gal}(E/K)\trianglelefteq G$ and $\mathrm{Gal}(K/F)\cong G/\mathrm{Gal}(E/K)$.
  \item $K = E_{\mathrm{Gal}(E/K)}$.
  \item $H\leq G \Rightarrow E_H$ is fixed by exactly $H$, i.e., $\mathrm{Gal}(E/E_H)=H$.
\end{enumerate}
\end{theorem}

\begin{proof}
\textbf{(1)} By induction on $[E:F]=n$.

Base: $n=1 \Rightarrow E=F$, trivial.

Inductive step: Let $[E:F]=n\geq 2$, $E$ a splitting field of $f(x)\in F[x]$.
Let $p(x)\mid f(x)$ be irreducible of degree $m=\deg(p)\geq 2$.
Since $\mathrm{char}(F)=0$ or $F$ finite, $p(x)$ has $m$ distinct roots $u=u_1,\ldots,u_m\in E$.

Let $G=\mathrm{Gal}(E/F)$, $R=\{u_1,\ldots,u_m\}$.
For each $u_i$, $\sigma\mapsto u_i$ is an algebraic conjugate of $p(x)$, so $R=\{\sigma(u)\mid\sigma\in G\}$.

Set $H=\{\sigma\in G\mid \sigma(u)=u\}$, $X=\{\sigma H\mid\sigma\in G\}$.
Then $H\leq G$, $|G|=[G:H]|H|=|X||H|$.

Note: $\sigma H=\tau H \Leftrightarrow \tau^{-1}\sigma\in H \Leftrightarrow \sigma(u)=\tau(u)$.
So $\phi: X\to R$, $\phi(\sigma H)=\sigma(u)$ is a bijection, giving $|X|=|R|=m$.

Hence $|G|=m|H|=[F(u):F]|H|$.

Since $\sigma\in\mathrm{Gal}(E/F(u))$ fixes $u$, we get $\mathrm{Gal}(E/F(u))\subseteq H$.
Conversely, $\sigma\in H\Rightarrow \sigma(u)=u$ and each $v\in F(u)$ writes as
$v=a_0+a_1u+\cdots+a_{m-1}u^{m-1}$, so $\sigma(v)=v$, giving $H\subseteq\mathrm{Gal}(E/F(u))$.

Thus $H=\mathrm{Gal}(E/F(u))$ and $|H|=|\mathrm{Gal}(E/F(u))|=[E:F(u)]$ (induction hypothesis, since $[E:F(u)]<n$).

Therefore $|G|=[F(u):F][E:F(u)]=[E:F]=n$, completing (1).

Also, $[K:F]=|G|/|\mathrm{Gal}(E/K)|$ follows since
$\sigma\mathrm{Gal}(E/K)=\tau\mathrm{Gal}(E/K) \Leftrightarrow \tau^{-1}\sigma\in\mathrm{Gal}(E/K)
\Leftrightarrow \sigma(\alpha)=\tau(\alpha)\ \forall\alpha\in K$,
so $[K:F]=|G|/|\mathrm{Gal}(E/K)|$.

\medskip
\textbf{(3) \& (4)} By definition $K\subseteq E_{\mathrm{Gal}(E/K)}$.
For $\alpha\in E\setminus K$: since $E/F$ is a splitting field, $\exists\,\sigma\in\mathrm{Gal}(E/K)$ with $\sigma(\alpha)\neq\alpha$,
so $\alpha\notin E_{\mathrm{Gal}(E/K)}$. Thus $E_{\mathrm{Gal}(E/K)}=K$.

For (4): $H\leq\mathrm{Gal}(E/E_H)$ by definition.
Suppose $|H|<|\mathrm{Gal}(E/E_H)|=n$; set $H=\{\sigma_1,\ldots,\sigma_{|H|}\}$ and $f(x)=\prod_{i=1}^{|H|}(x-\sigma_i(\alpha))$ for some $\alpha\in K$.
Coefficients of $f$ are symmetric in $\sigma_i(\alpha)$'s, hence in $E_H$.
But $f(\alpha)=0$ gives $\deg(\alpha,E_H)\leq|H|<n=[E:E_H]$, contradiction.
So $H=\mathrm{Gal}(E/E_H)$.

\medskip
\textbf{(2)} $K/F$ Galois $\Rightarrow$ $K$ is a splitting field of some $g(x)\in F[x]$.
By (3), $K=E_{\mathrm{Gal}(E/K)}$, and for $\sigma\in\mathrm{Gal}(E/F)$, $\alpha\in K$:
$\sigma(\alpha)\in K$, so $\tau(\sigma(\alpha))=\sigma(\alpha)$ for all $\tau\in\mathrm{Gal}(E/K)$,
i.e., $(\sigma^{-1}\tau\sigma)(\alpha)=\alpha$, giving $\mathrm{Gal}(E/K)\trianglelefteq G$.

Define $\phi: G\to\mathrm{Gal}(K/F)$ by $\phi(\sigma)=\sigma_K$ (restriction).
$\ker(\phi)=\mathrm{Gal}(E/K)$, so by the first isomorphism theorem,
\[
  \mathrm{Gal}(K/F)\cong G/\mathrm{Gal}(E/K). \qedhere
\]
\end{proof}

\end{document}










\end{document}



\end{document}




\end{document}






\end{document}
